% ===== Wrap-up (~5 min, then questions) =====
\partaccent{butter}\notopic

\begin{frame}{Այսօրվա 5 գաղափարը}
  \begin{enumerate}
    \item<1-> Ուղեղը ապշեցուցիչ է, բայց \textbf{կատարյալ չէ}. էվոլյուցիան նորը կառուցել է հնի վրայից
    \item<2-> Հին ու նոր ուղեղները անընդհատ վիճում են. երբեմն հաղթում է մեկը (ծովախտ), երբեմն՝ մյուսը (աղանդեր)
    \item<3-> Քունը պարտադիր է, թեև դեռ չգիտենք՝ ինչու։ Դեռահասին պետք է \textbf{8-10 ժամ}
    \item<4-> Վախը նախ կրակում է, հետո մտածում. ավելի լավ է ապահով, քան ափսոս
    \item<5-> Հիշողությունը ձայնագրություն չէ. այն ամեն անգամ \textbf{վերակառուցվում} ու մի քիչ խմբագրվում է
  \end{enumerate}
\end{frame}

\begin{frame}{Փորձիր տանը}
  \begin{columns}[T]
    \begin{column}{0.31\textwidth}
      \begin{card}[colback=mint!30!cardwhite]
        \centering\bigemoji[1cm]{potato}\ \bigemoji[1cm]{apple}\par\smallskip
        \small\raggedright Կարտոֆիլ թե խնձոր. փակ աչքերով ու քթով ստուգիր ընտանիքիդ անդամներին
      \end{card}
    \end{column}
    \begin{column}{0.31\textwidth}
      \begin{card}[colback=lavender!35!cardwhite]
        \centering\bigemoji[1cm]{wavehand}\par\smallskip
        \small\raggedright Նոր ծանոթի անունը կրկնիր ու զրույցում երկու անգամ օգտագործիր։ Վաղը կհիշե՞ս
      \end{card}
    \end{column}
    \begin{column}{0.31\textwidth}
      \begin{card}[colback=peach!35!cardwhite]
        \centering\bigemoji[1cm]{bed}\par\smallskip
        \small\raggedright Մեկ շաբաթ գրանցիր, թե քանի ժամ ես քնում։ Հասնու՞մ ես 8-ի
      \end{card}
    \end{column}
  \end{columns}
\end{frame}

\begin{frame}{Գրքի հաջորդ գլուխները}
  \begin{columns}[c]
    \begin{column}{0.6\textwidth}
      \chip[brainpink]{scream}{3. Վախ. վախենալու բան չկա}\par\smallskip
      {\small Սնահավատություն ու դավադրության տեսություններ, ֆոբիաներ, ինչու ենք սիրում վախենալ, ու ինչու է քննադատությունն ավելի ուժեղ, քան գովեստը}\par\bigskip
      \chip[butter]{bulb}{4. Կարծում ես՝ խելացի՞ ես}\par\smallskip
      {\small Ինչ է չափում IQ-ն, ինչու են խելացի մարդիկ հաճախ պարտվում վեճերում, ու օգնու՞մ են արդյոք խաչբառերը}
    \end{column}
    \begin{column}{0.34\textwidth}
      \centering\mascot[3cm]{scared}
    \end{column}
  \end{columns}
\end{frame}

\begin{frame}[plain]
  \begin{tikzpicture}[remember picture, overlay]
    \fill[lavender!60!cream] (current page.south west) rectangle (current page.north east);
    \fill[pattern={Dots[distance=4.4mm, radius=0.28mm]}, pattern color=plum!10!lavender!60!cream]
      (current page.south west) rectangle (current page.north east);
    \node at ([xshift=-3.2cm]current page.center) {\mascot[4.4cm]{happy}};
    \node[anchor=west, text=plum] at ([xshift=0.2cm, yshift=0.6cm]current page.center)
      {\titlefont\fontsize{40}{42}\selectfont Հարցեր};
    \node[anchor=west, text=plum, font=\mediumfont\large] at ([xshift=0.2cm, yshift=-0.8cm]current page.center)
      {Գրեք չաթում կամ միացրեք խոսափողը \emoji{megaphone}};
    \node[anchor=west, text=plum!70!cream, font=\small] at ([xshift=0.2cm, yshift=-1.6cm]current page.center)
      {Դին Բըրնեթ, «The Idiot Brain», 2016};
  \end{tikzpicture}
\end{frame}
